\chapter{Örnekleme Yöntemleri ve Örnekleme Yanlılığı}
\label{ch:sampling-design}
\index{örnekleme}
\index{örnekleme yanlılığı}

\begin{hedefler}
Bu bölümün sonunda okuyucu:
\begin{itemize}
  \item hedef evren, erişilebilir evren, örnekleme çerçevesi, örnekleme birimi ve örneklemi ayırt edebilecek,
  \item basit rastgele, sistematik, tabakalı, küme ve çok aşamalı örneklemeyi araştırma bağlamına göre seçebilecek,
  \item örnekleme hatası ile kapsam, yanıtlamama, gönüllülük ve ölçüm yanlılığını ayırabilecek,
  \item örneklem hacmini hassasiyet, güç, tasarım ve beklenen yanıt oranıyla ilişkilendirebilecek,
  \item örnekleme sürecini şeffaf ve etik biçimde raporlayabilecektir.
\end{itemize}
\end{hedefler}

\begin{hazirlik}
Bir anket yalnız sosyal medya hesabını takip eden öğrencilere gönderilirse
kimler sistematik olarak dışarıda kalabilir? Genellenebilirliği tartışın.
\end{hazirlik}

\section{PDR vakası: Üniversite öğrencilerinin yardım arama tutumları}

Bir üniversitenin psikolojik danışma merkezi, öğrencilerin profesyonel yardım
arama tutumlarını ve merkezin hizmetlerine ilişkin farkındalıklarını
incelemek istemektedir. Üniversitede dört fakültede toplam 6.000 öğrenci
vardır. Araştırma ekibinin bütün öğrencilere ulaşması mümkün değildir; bu
nedenle 600 öğrencilik bir örneklem planlanmaktadır.

İlk bakışta yapılacak iş 600 isim seçmek gibi görünür. Oysa şu kararlar
sonucun niteliğini belirler:

\begin{itemize}
  \item Araştırmanın hedefi bütün kayıtlı öğrenciler mi, yalnız kampüse devam edenler mi?
  \item Güncel öğrenci listesi bütün fakülte ve sınıf düzeylerini içeriyor mu?
  \item Az sayıdaki alt gruplar örneklemde yeterince temsil edilecek mi?
  \item Ankete yanıt vermeyen öğrenciler yanıt verenlerden farklı olabilir mi?
  \item Yardım arama gibi hassas bir konuda davet biçimi gönüllülüğü etkiler mi?
\end{itemize}

Bu vaka bölüm boyunca farklı örnekleme tasarımlarını karşılaştırmak için
kullanılacaktır.

\section{Hedef evrenden gözlenen örnekleme}
\index{hedef evren}
\index{erişilebilir evren}
\index{örnekleme çerçevesi}
\index{örnekleme birimi}
\index{analiz birimi}

Araştırmacının hakkında sonuç üretmek istediği topluluk \textbf{hedef evren},
seçim yapılırken fiilen erişilebilen birimlerin listesi \textbf{örnekleme
çerçevesi}, ölçülen birimler ise \textbf{örneklem}dir. Bu üç küme aynı değilse
genelleme yapılırken kapsam hatası oluşabilir \cite{cochran1977,lohr2021}.

\begin{description}
  \item[Erişilebilir evren] Araştırmacının uygulamada ulaşabildiği hedef evren bölümüdür.
  \item[Örnekleme birimi] Seçim işleminin uygulandığı öğrenci, sınıf, okul veya kurumdur.
  \item[Analiz birimi] Sonuçların hesaplandığı temel birimdir; örnekleme birimiyle aynı olmak zorunda değildir.
\end{description}

Örneğin önce sınıflar seçilip seçilen sınıflardaki öğrenciler ölçülürse sınıf
örnekleme birimi, öğrenci analiz birimidir. Bu ayrım bağımsızlık ve standart
hata hesapları açısından önemlidir.

Örneğin bir üniversitedeki bütün öğrenciler hedeflenirken anket bağlantısının
yalnız belirli bir dersin çevrim içi sayfasında paylaşılması, örnekleme
çerçevesini o dersi alan öğrencilerle sınırlar.

\begin{figure}[htbp]
\centering
\resizebox{0.92\linewidth}{!}{%
\begin{tikzpicture}[alt={Hedef evren, erişilebilir evren, örnekleme çerçevesi ve gözlenen örneklem arasındaki daralmayı gösteren akış şeması},
  box/.style={draw=PrismBlue,rounded corners,fill=PrismLight,minimum width=3.2cm,minimum height=1cm,align=center},
  note/.style={font=\scriptsize,text=PrismGray,align=center}
]
\node[box] (target) at (0,0) {Hedef evren\\6.000 öğrenci};
\node[box] (frame) at (4.3,0) {Örnekleme çerçevesi\\güncel öğrenci listesi};
\node[box] (selected) at (8.6,0) {Seçilen örneklem\\600 öğrenci};
\node[box] (responded) at (12.9,0) {Yanıtlayan örneklem\\örneğin 438 öğrenci};
\draw[->,thick,PrismBlue] (target)--(frame);
\draw[->,thick,PrismBlue] (frame)--(selected);
\draw[->,thick,PrismBlue] (selected)--(responded);
\node[note] at (2.15,-0.85) {kapsam hatası};
\node[note] at (6.45,-0.85) {örnekleme hatası};
\node[note] at (10.75,-0.85) {yanıtlamama};
\end{tikzpicture}%
}
\caption{Hedef evrenden yanıtlayan örnekleme uzanan seçim zinciri.}
\end{figure}

\subsection{İyi bir örnekleme çerçevesinin özellikleri}
\index{örnekleme çerçevesi!kapsam}
\index{kapsam hatası}

Çerçeve güncel olmalı, hedef evrendeki her uygun birimi kapsamalı, aynı birimi
birden fazla kez içermemeli ve seçim için gerekli tabaka ya da küme bilgisini
taşımalıdır. Mezun olmuş öğrencilerin listede kalması fazla kapsam; uzaktan
öğrenim öğrencilerinin listede bulunmaması eksik kapsam örneğidir.

\begin{etkinlik}
Yardım arama tutumu araştırması için üç olası çerçeve düşünün: öğrenci bilgi
sistemi listesi, kampüste o gün bulunan öğrenciler ve merkezin sosyal medya
takipçileri. Her çerçevenin kimleri kapsadığını ve kimleri dışarıda
bıraktığını karşılaştırın.
\end{etkinlik}

\section{Olasılıklı örnekleme yöntemleri}
\index{olasılıklı örnekleme}
\index{dahil olma olasılığı}

Olasılıklı örneklemede her birimin seçilme olasılığı bilinir ve sıfırdan
büyüktür. Bu özellik, örnekleme belirsizliğinin hesaplanmasına ve tasarıma
dayalı genellemeye temel sağlar. ``Rastgele'' sözcüğü gelişigüzel seçim değil,
önceden tanımlanmış bir olasılık mekanizması anlamına gelir.

\begin{description}
  \item[Basit rastgele örnekleme] Her birim bilinen ve eşit seçilme şansına sahiptir.
  \item[Sistematik örnekleme] Rastgele başlangıçtan sonra listedeki her $k$'ıncı birim seçilir.
  \item[Tabakalı örnekleme] Evren önemli alt gruplara ayrılır ve her tabakadan örnek alınır.
  \item[Küme örneklemesi] Bireyler yerine sınıf, okul veya kurum gibi doğal kümeler seçilir.
  \item[Çok aşamalı örnekleme] Önce kümeler, sonra kümeler içindeki alt birimler seçilir.
\end{description}

\subsection{Basit rastgele örnekleme}
\index{basit rastgele örnekleme}

$N$ birimlik çerçeveden $n$ birim seçildiğinde her birimin dahil olma
olasılığı
\[
\pi_i=\frac{n}{N}
\]
olur. Seçim bilgisayarın rassal sayı üreteciyle yapılabilir; araştırmacının
``rastgele görünen'' isimler seçmesi yeterli değildir.

\begin{ornek}[Basit rastgele seçim]
6.000 öğrenciden 600'ü seçilecekse her öğrencinin örnekleme dahil olma
olasılığı $600/6000=0{,}10$'dur. Listenin her satırı tek ve uygun bir öğrenciyi
temsil ediyorsa 600 farklı sıra numarası rassal olarak seçilebilir.
\end{ornek}

\subsection{Sistematik örnekleme}
\index{sistematik örnekleme}
\index{seçim aralığı}

Yaklaşık seçim aralığı $k=N/n$ ile bulunur. $N=6000$ ve $n=600$ için
$k=10$'dur. 1 ile 10 arasında rassal bir başlangıç, örneğin 7 seçilir; sonra
7, 17, 27, $\ldots$ numaralı birimler alınır.

Liste sırası gizli bir döngü içeriyorsa sistematik seçim yanlı sonuç
üretebilir. Örneğin kayıtlar her on satırda bir aynı program türünü tekrar
ediyorsa $k=10$ seçimi bazı grupları aşırı ya da eksik temsil edebilir.

\subsection{Tabakalı örnekleme}
\index{tabakalı örnekleme}
\index{orantılı dağıtım}

Evren fakülte, sınıf düzeyi veya öğrenim türü gibi önemli alt gruplara
ayrılır. Orantılı dağıtımda $h$ tabakasından seçilecek sayı
\[
n_h=n\frac{N_h}{N}
\]
olarak belirlenir. Fakülte büyüklükleri 2.400, 1.800, 1.200 ve 600 ise
600 kişilik orantılı örneklem sırasıyla 240, 180, 120 ve 60 öğrenciden oluşur.

Küçük fakat araştırma açısından önemli bir alt grup orantılı seçimde çok az
gözlemle temsil edilebilir. Bu durumda orantısız tabakalama yapılabilir; ancak
üniversite geneli sonuçlarda seçim olasılıklarını dengeleyen ağırlıklar
gerekebilir.

\subsection{Küme ve çok aşamalı örnekleme}
\index{küme örneklemesi}
\index{çok aşamalı örnekleme}
\index{tasarım etkisi}

Öğrencilerin tek listesine ulaşmak güç, sınıf listelerine ulaşmak kolaysa önce
sınıflar seçilip sonra seçilen sınıflardaki öğrenciler araştırmaya alınabilir.
Bu yaklaşım maliyeti azaltır. Bununla birlikte aynı sınıftaki öğrenciler
birbirine benzeyebileceğinden 600 öğrencilik küme örneklemi, 600 bağımsız
öğrencilik basit rastgele örneklem kadar bilgi taşımayabilir.

Çok aşamalı tasarımda önce fakülteler veya sınıflar, sonra bunların içindeki
öğrenciler seçilir. Her aşamadaki seçim olasılığı ve küme yapısı analizde
dikkate alınmalıdır.

Tabakalama, alt grupların temsilini güvenceye almak için; küme örneklemesi ise
dağınık bir evrende veri toplama maliyetini azaltmak için kullanılabilir.
Küme içindeki bireyler birbirine benzediğinde etkili bilgi miktarı, aynı
sayıdaki bağımsız bireyden daha düşük olur \cite{chambers2003}.

\begin{table}[htbp]
\centering
\caption{Olasılıklı örnekleme yöntemlerinin karşılaştırılması.}
\begin{tabular}{@{}p{0.20\textwidth}p{0.27\textwidth}p{0.29\textwidth}@{}}
\toprule
Yöntem & Güçlü yön & Temel risk veya maliyet \\
\midrule
Basit rastgele & Seçim mantığı ve ağırlıklar yalındır & Tam ve güncel birey listesi gerekir \\
Sistematik & Uygulaması hızlı ve örneklemi listeye yayar & Liste sırasındaki dönemlilik sonucu bozabilir \\
Tabakalı & Önemli alt grupların temsilini güvenceye alır & Tabaka bilgisi ve gerektiğinde ağırlık gerekir \\
Küme & Saha maliyetini azaltır & Küme içi benzerlik hassasiyeti düşürebilir \\
Çok aşamalı & Büyük ve dağınık evrenlerde uygulanabilirdir & Seçim ve analiz yapısı daha karmaşıktır \\
\bottomrule
\end{tabular}
\end{table}

\section{Olasılıksız örnekleme}
\index{olasılıksız örnekleme}
\index{kolayda örnekleme}
\index{amaçlı örnekleme}
\index{kartopu örnekleme}

Kolayda, gönüllü ve amaçlı örnekleme bazı araştırma bağlamlarında gerekli
olabilir; ancak birimlerin seçilme olasılıkları bilinmediği için örnekleme
hatası klasik formüllerle tam olarak hesaplanamaz. Sonuçların hangi gruba
genellenebileceği açıkça sınırlandırılmalıdır.

\begin{description}
  \item[Kolayda örnekleme] Ulaşılması kolay katılımcılar seçilir.
  \item[Gönüllü örnekleme] Katılımcılar açık çağrıya kendi istekleriyle yanıt verir.
  \item[Amaçlı örnekleme] Belirli deneyim veya özellikleri taşıyan bilgi-zengin durumlar seçilir.
  \item[Kartopu örnekleme] Ulaşılması güç gruplarda katılımcılar başka katılımcılara erişim sağlar.
  \item[Kota örneklemesi] Belirli grup sayıları hedeflenir; grup içi seçim olasılıklı olmak zorunda değildir.
\end{description}

Amaçlı veya kartopu örnekleme nitel ve keşfedici PDR araştırmalarında değerli
olabilir. Sorun yöntemin kullanılması değil, sonuçların olasılıklı bir örneklem
elde edilmiş gibi bütün evrene genellenmesidir.

\begin{researchbox}
Klinik deneyimi bulunan az sayıdaki danışmanın görüşünü derinlemesine
incelemek için amaçlı örnekleme uygun olabilir. Aynı örneklemden ülkedeki tüm
danışmanların belirli bir görüşü hangi oranda taşıdığını tahmin etmek ise
savunulamaz.
\end{researchbox}

\section{Yanlılık kaynakları}
\index{yanlılık!kapsam}
\index{yanlılık!gönüllülük}
\index{yanlılık!yanıtlamama}
\index{yanlılık!ölçüm}

\begin{itemize}
  \item \textbf{Kapsam yanlılığı:} Çerçeve hedef evrenin bazı kesimlerini dışlar.
  \item \textbf{Yanıtlamama yanlılığı:} Yanıt verenler vermeyenlerden sistematik olarak farklıdır.
  \item \textbf{Gönüllülük yanlılığı:} Güçlü görüşü olan kişiler katılmaya daha isteklidir.
  \item \textbf{Ölçüm yanlılığı:} Soru biçimi, araç veya uygulayıcı yanıtları sistematik etkiler.
  \item \textbf{Seçim sonrası dışlama:} Analiz ölçütleri belirli sonuçları aşırı temsil eder.
\end{itemize}

\subsection{Örnekleme hatası ve yanlılık aynı değildir}
\index{örnekleme hatası}
\index{yanlılık}

Örnekleme hatası, yalnız bir örneklemin gözlenmesinden kaynaklanan rassal
farktır; farklı örneklemler farklı sonuç verir. Örneklem büyüdükçe çoğunlukla
azalır ve standart hatayla nicelendirilebilir. Yanlılık ise sonuçları belirli
bir yönde sistematik olarak çarpıtır; daha büyük örneklem bu çarpıklığı
gidermeyebilir.

\begin{table}[htbp]
\centering
\caption{Örnekleme sürecindeki hata ve yanlılıkların tanılanması.}
\begin{tabular}{@{}p{0.22\textwidth}p{0.30\textwidth}p{0.29\textwidth}@{}}
\toprule
Sorun & Yardım arama araştırmasındaki örnek & Olası önlem \\
\midrule
Kapsam & Kayıt donduran öğrenciler listede yok & Hedef evreni sınırla veya çerçeveyi tamamla \\
Yanıtlamama & Hizmetlere ilgisiz öğrenciler ankete dönmüyor & Hatırlatma, kısa form, yanıtlayan--yanıtlamayan karşılaştırması \\
Gönüllülük & Olumlu veya olumsuz deneyimi güçlü olanlar katılıyor & Olasılıklı davet ve tarafsız çağrı dili \\
Ölçüm & ``Merkezimizin yararlı hizmetleri'' biçiminde yönlendirici soru & Tarafsız ifade ve pilot uygulama \\
İşleme & Eksik yanıtların gerekçesiz silinmesi & Önceden tanımlı veri temizleme ve duyarlılık analizi \\
\bottomrule
\end{tabular}
\end{table}

\begin{mistakebox}
Rastgele örnekleme ile rastgele atama aynı işlem değildir. Rastgele örnekleme
evrene genellemeyi, rastgele atama ise gruplar arasında nedensel karşılaştırmayı
destekler.
\end{mistakebox}

\section{Yanıtlamama ve katılım süreci}
\index{yanıtlamama}
\index{yanıt oranı}

Yanıt oranı genel olarak
\[
\text{yanıt oranı}=\frac{\text{uygun ve tamamlanmış yanıt sayısı}}
{\text{ulaşılabilen uygun davetli sayısı}}
\]
biçiminde düşünülür. Ancak yüksek yanıt oranı tek başına yanlılığın olmadığını,
düşük oran da mutlaka ağır yanlılık bulunduğunu kanıtlamaz. Önemli olan
yanıtlamanın araştırma değişkenleriyle sistematik ilişkili olup olmadığıdır.

Katılımı artırmak için açık ve kısa davet, makul süre, erişilebilir anket,
birkaç tarafsız hatırlatma ve gizlilik güvencesi kullanılabilir. Ders notu,
hizmet erişimi veya danışmanlık ilişkisi üzerinden baskı oluşturacak teşvik ve
çağrılardan kaçınılmalıdır.

\begin{etkinlik}
Aynı ankete ait iki davet metni yazın. Birincisi öğrencileri katılmaya
yönlendiren ve araştırmanın beklenen sonucunu ima eden sorunlu bir metin,
ikincisi tarafsızlık, gönüllülük ve gizlilik ilkelerini koruyan etik bir metin
olsun. Metinleri karşılaştırın.
\end{etkinlik}

\section{Örneklem hacmi ve temsil}
\index{örneklem hacmi}
\index{temsil gücü}
\index{sonlu evren düzeltmesi}

Büyük örneklem rastgele örnekleme hatasını azaltır; fakat sistematik yanlılığı
otomatik olarak gidermez. Örneklem hacmi belirlenirken hedeflenen hassasiyet,
değişkenlik, güven düzeyi, araştırma tasarımı ve ulaşılabilir kaynaklar birlikte
değerlendirilir.

Basit rastgele örneklemede ortalamanın standart hatası yaklaşık
\[
SE(\bar X)=\frac{s}{\sqrt n}
\]
olduğundan belirsizliği yarıya indirmek yaklaşık dört kat gözlem gerektirir.

Örneklem hacmi için tek bir evrensel ``yeterli sayı'' yoktur. Planlama şu
bileşenlere dayanır:

\begin{itemize}
  \item tahmin edilmek istenen etki veya oran,
  \item kabul edilebilir hata payı ve güven düzeyi,
  \item evrendeki değişkenlik,
  \item karşılaştırılacak grup ve alt grup sayısı,
  \item kümeleme veya ağırlıklandırma gibi tasarım özellikleri,
  \item beklenen yanıtlamama ve veri kaybı,
  \item zaman, bütçe ve etik yük.
\end{itemize}

Örneğin tamamlanmış 600 yanıt hedefleniyor ve uygun davetlilerin yaklaşık
yüzde 75'inin yanıt vereceği bekleniyorsa en az
\[
600/0{,}75=800
\]
öğrencinin davet edilmesi gerekir. Bu hesap yanıtlamama yanlılığını çözmez;
yalnız beklenen kayıp için sayısal planlama yapar.

\begin{sezgibox}
Örneklem hacmi hassasiyeti, örnekleme tasarımı ise temsilin dayanağını
belirler. Çok büyük ama seçimi yanlı bir örneklem, küçük standart hata ile
yanlış değeri çok hassas biçimde tahmin edebilir.
\end{sezgibox}

\section{Adım adım çözümlü örnekleme planı}

Yardım arama araştırmasında fakülte büyüklükleri sırasıyla 2.400, 1.800,
1.200 ve 600; hedef tamamlanmış örneklem 600 olsun.

\subsection*{1. Tasarım seçimi}

Yardım arama tutumunun fakültelere göre değişebileceği ve her fakültenin
örneklemde görünmesi istendiği için fakülte tabakalı örnekleme seçilir. Her
tabaka içinde öğrenciler basit rastgele yöntemle belirlenir.

\subsection*{2. Orantılı dağıtım}

Fakültelere ayrılacak tamamlanmış örneklem sayıları
\[
600\left(\frac{2400}{6000}\right)=240,\quad
600\left(\frac{1800}{6000}\right)=180,\quad
600\left(\frac{1200}{6000}\right)=120,\quad
600\left(\frac{600}{6000}\right)=60
\]
olarak bulunur.

\subsection*{3. Yanıtlamama payı}

Her fakültede yüzde 75 yanıt beklenirse davet sayıları yaklaşık 320, 240, 160
ve 80; toplam 800 olur. Gerçekleşen yanıt oranları fakülteler arasında
karşılaştırılmalı ve dengesizlik raporlanmalıdır.

\subsection*{4. Seçim ve kayıt}

Her öğrencinin çerçevedeki benzersiz kimliği, tabakası, seçilip seçilmediği,
davet tarihi ve yanıt durumu güvenli bir seçim günlüğünde tutulur. Analiz
dosyasındaki kimlik bilgileri bu yönetim dosyasından ayrılır.

\section{Python ile tabakalı örnek seçimi}

\begin{lstlisting}
import pandas as pd

# cerceve: her satir bir ogrenci; id ve fakulte sutunlari var
cerceve = pd.read_csv("ogrenci_cercevesi.csv")
davet_sayilari = {"A": 320, "B": 240, "C": 160, "D": 80}

parcalar = []
for fakulte, n in davet_sayilari.items():
    tabaka = cerceve.loc[cerceve["fakulte"] == fakulte]
    parcalar.append(tabaka.sample(n=n, replace=False, random_state=209))

orneklem = pd.concat(parcalar).sort_values(["fakulte", "id"])
print(orneklem["fakulte"].value_counts().sort_index())
orneklem.to_csv("davet_orneklemi.csv", index=False)
\end{lstlisting}

Rassal tohumun kaydedilmesi seçimin yeniden üretilebilmesini sağlar. Gerçek
bir projede aynı sabit tohum her tabakada aynen kullanılmak yerine tek bir
rassal sayı üreteci yönetilebilir; burada kod öğretim amacıyla sade
tutulmuştur.

\section{SPSS ile rastgele seçim}

Tam bir çerçeve SPSS veri dosyasında bulunuyorsa \textbf{Data $>$ Select Cases
$>$ Random sample of cases} menüsü basit rastgele alt örneklem üretmek için
kullanılabilir. Tabakalı seçimde önce fakülteye göre dosya bölme veya her
tabaka için ayrı seçim gerekir. Seçilmeyen kayıtların silinmesi yerine filtre
değişkeni kullanmak, seçim işlemini denetlenebilir tutar.

\begin{outputguidebox}
Seçimden sonra yalnız toplam $n$ değerini kontrol etmeyin. Her tabakadaki
seçilen sayı, yinelenen kimlikler, kapsam dışı kayıtlar ve beklenen seçim
oranları incelenmelidir. Rastgele seçim komutu kötü bir örnekleme çerçevesini
düzeltemez.
\end{outputguidebox}

\section{Örnekleme sürecini raporlama}

Yeterli bir yöntem bölümü hedef evreni, çalışma tarihini, uygunluk ve dışlama
ölçütlerini, çerçevenin kaynağını, örnekleme yöntemini, planlanan ve gerçekleşen
tabaka sayılarını, davet ve yanıt sayılarını, yanıt oranını ve varsa ağırlıkları
belirtir. Örnek rapor:

``Araştırmanın hedef evrenini 2026--2027 akademik yılında üniversiteye kayıtlı
6.000 öğrenci oluşturmuştur. Güncel öğrenci bilgi sistemi listesi fakülteye
göre tabakalanmış; her tabakada bilgisayar üretimli rassal sayılarla toplam
800 öğrenci davet edilmiştir. 612 öğrenci anketi tamamlamıştır (yanıt oranı
\%76{,}5). Fakülteye göre planlanan ve gerçekleşen yanıt sayıları raporlanmış,
yanıtlamama nedeniyle oluşabilecek yanlılık araştırmanın sınırlılığı olarak
değerlendirilmiştir.''

\section{Literatür veri setiyle IBM SPSS laboratuvarı}
\index{Student Performance veri seti!tabakalı örnekleme}
\index{SPSS!tabakalı örnekleme}
\index{rassal tohum}
\index{dahil olma olasılığı}

Birinci bölümde içe aktarılan \emph{Student Performance} matematik dosyası bu
uygulamada 395 birimlik sonlu bir örnekleme çerçevesi gibi ele alınacaktır
\cite{cortez2008data,cortezsilva2008}. Amaç, iki okulun her ikisini de temsil
eden yaklaşık yüzde 25 büyüklüğünde bir öğretim örneklemi seçmektir. Çerçevede
Gabriel Pereira (GP) okulundan 349, Mousinho da Silveira (MS) okulundan 46
öğrenci vardır. Orantılı dağıtım sonucunda 100 kişilik örnekleme GP'den 88,
MS'den 12 öğrenci alınacaktır.

\begin{researchbox}
Buradaki işlem yayımlanmış veri dosyasından yeniden örnek seçme alıştırmasıdır.
Özgün araştırmadaki öğrencileri sonradan olasılıklı bir evren örneklemine
dönüştürmez. Bu nedenle elde edilen sonuçlar Portekiz'deki bütün öğrencilere
tasarıma dayalı olarak genellenemez.
\end{researchbox}

\subsection{Araştırma görevi ve seçim planı}

\begin{enumerate}
  \item Çerçevedeki okul frekanslarını doğrulayın.
  \item Her okul içinde rassal sayı üretip öğrencileri sıralayın.
  \item GP tabakasından ilk 88, MS tabakasından ilk 12 öğrenciyi seçin.
  \item Seçilen sayıları ve tabaka bazındaki dahil olma oranlarını denetleyin.
  \item Seçim göstergesini koruyun; seçilmeyen kayıtları veri dosyasından silmeyin.
\end{enumerate}


Sabit rassal tohum aynı yazılım ve aynı veri sırası altında seçimin yeniden
üretilmesini kolaylaştırır. Tohum, seçimin araştırmacı tarafından istenen
sonucu verecek biçimde tekrar tekrar denenmesi için kullanılmamalıdır.

\subsection{SPSS syntax ile orantılı tabakalı seçim}

Aşağıdaki syntax, birinci bölümde oluşturulan \texttt{StudentMath} veri seti
açıkken çalıştırılır. Önce özgün satır sırasını saklayan bir kimlik ve okul
tabakalarının büyüklükleri oluşturulur; ardından her tabaka içinde rassal
sıralama yapılır.

\begin{lstlisting}[language={}]
DATASET ACTIVATE StudentMath.

COMPUTE case_id=$CASENUM.
AGGREGATE
 /OUTFILE=* MODE=ADDVARIABLES
 /BREAK=school
 /N_tabaka=N.

SET SEED=2092026.
COMPUTE rassal=RV.UNIFORM(0,1).
SORT CASES BY school rassal.
RANK VARIABLES=rassal (A) BY school
 /RANK INTO tabaka_sira.

DO IF school='GP'.
 COMPUTE hedef_n=88.
ELSE IF school='MS'.
 COMPUTE hedef_n=12.
END IF.

COMPUTE secildi=(tabaka_sira<=hedef_n).
COMPUTE dahil_olma=hedef_n/N_tabaka.
VALUE LABELS secildi 0 'Seçilmedi' 1 'Seçildi'.
FORMATS dahil_olma (F5.3).

CROSSTABS
 /TABLES=school BY secildi
 /CELLS=COUNT ROW.

TEMPORARY.
SELECT IF secildi=1.
FREQUENCIES VARIABLES=school.

TEMPORARY.
SELECT IF secildi=1.
DESCRIPTIVES VARIABLES=G3 absences
 /STATISTICS=MEAN STDDEV MIN MAX.

TEMPORARY.
SELECT IF secildi=1.
SAVE OUTFILE='student-mat-tabakali-100.sav'.

SORT CASES BY case_id.
EXECUTE.
\end{lstlisting}

\subsection{SPSS çıktısının erişilebilir özeti}

Seçimin temel denetimi, \textbf{Crosstabs} tablosundaki gözlem sayıları ve satır
yüzdeleriyle yapılır:

\begin{table}[htbp]
\centering
\caption{Okula göre orantılı tabakalı seçim denetimi.}
\begin{tabular}{@{}lrrrr@{}}
\toprule
Okul & Çerçeve $N_h$ & Seçilen $n_h$ & Seçilmeyen & Dahil olma oranı \\
\midrule
GP & 349 & 88 & 261 & 0,252 \\
MS & 46  & 12 & 34  & 0,261 \\
\midrule
Toplam & 395 & 100 & 295 & 0,253 \\
\bottomrule
\end{tabular}
\end{table}

\subsection{Çıktının analizi}

Toplam seçilen gözlem sayısı planlanan 100'dür ve iki okul da örneklemde
temsil edilmektedir. Tam sayıya yuvarlama nedeniyle dahil olma oranları tam
olarak eşit değildir: GP için $88/349=0{,}252$, MS için
$12/46=0{,}261$'dir. Fark küçük olsa da kesin tasarım ağırlıkları kullanılacaksa
her öğrenciye kendi dahil olma olasılığının tersi verilmelidir:
\[
w_{GP}=\frac{349}{88}\approx3{,}97,
\qquad
w_{MS}=\frac{46}{12}\approx3{,}83.
\]

\textbf{Descriptives} tablosundaki G3 ve devamsızlık özetleri seçilmiş 100
öğrenciyi betimler; çerçevedeki 395 öğrencinin kesin özellikleri değildir.
Başka bir rassal tohum farklı öğrencileri ve dolayısıyla biraz farklı bir
ortalama ile standart sapmayı seçebilir. Bu değişim bir yazılım hatası değil,
örnekleme değişkenliğidir.

Çapraz tabloda bir okuldan 88 ve diğerinden 12 yerine farklı sayılar
görülüyorsa önce filtre, okul kodları, veri sırası ve \texttt{BY school}
ifadesi denetlenmelidir. Yalnız toplamın 100 olması yeterli değildir; örnekleme
planının tabaka bazında da gerçekleşmesi gerekir.

\begin{outputguidebox}
Bu uygulamada ana çıktı bir anlamlılık testi değildir. Doğru okuma sırası:
(1) çerçeve büyüklükleri, (2) hedef tabaka sayıları, (3) seçilen gerçek
sayılar, (4) dahil olma oranları, (5) filtre durumudur. Düşük bir $p$-değeri
kötü kurulmuş bir örnekleme çerçevesini düzeltemez.
\end{outputguidebox}

\section{R ile çözümlü uygulama}
\label{sec:r-b06}

\textbf{Soru.} Her sınıfta dört öğrenci bulunan 12 kişilik çerçeveden altı
kişilik basit rastgele örneklem ve her sınıftan ikişer kişilik tabakalı
örneklem seçin. Seçilme olasılıklarını ve ağırlıkları bulun.

\begin{lstlisting}[style=prismR]
set.seed(2026)
ogrenciler <- data.frame(
  id = 1:12, sinif = factor(rep(1:3, each = 4))
)
basit <- ogrenciler[sample.int(nrow(ogrenciler), 6), ]
indisler <- split(seq_len(nrow(ogrenciler)),
                  ogrenciler$sinif)
secilen <- unlist(lapply(indisler,
                        function(indis) sample(indis, 2)),
                  use.names = FALSE)
tabakali <- ogrenciler[secilen, ]
tabakali$secim_olasiligi <- 2 / 4
tabakali$agirlik <- 1 / tabakali$secim_olasiligi
print(basit)
print(tabakali)
table(tabakali$sinif)
sum(tabakali$agirlik)
stopifnot(nrow(basit) == 6,
          anyDuplicated(basit$id) == 0,
          all(table(tabakali$sinif) == 2))
\end{lstlisting}

\begin{outputguidebox}
\textbf{Çözüm ve kontrol değerleri.} İki tasarımda da altı farklı öğrenci
seçilir. Tabakalı seçimde sınıf frekansları kesin olarak $(2,2,2)$'dir;
basit rastgele seçim bunu garanti etmez. Her öğrencinin seçilme olasılığı
$6/12=2/4=0{,}5$, tasarım ağırlığı 2'dir. Altı ağırlığın toplamı 12'dir.
\end{outputguidebox}

\textbf{Yorum.} Bu örnekte tabakalar eşit büyüklüktedir ve eşit sayıda
birim seçilir. Eşit olmayan tabaka hacimlerinde her tabaka için $N_h/n_h$
ayrı hesaplanmalıdır. Ağırlıklı ortalama hesaplamak, tasarıma uygun standart
hata hesabını kendiliğinden sağlamaz.

\textbf{Pekiştirme ve yanıt.} Bir sınıfın çerçeveden tamamen çıkarılması
kapsam hatası yaratır; kalan öğrencileri rastgele seçmek bu hatayı düzeltmez.
\section{Bölüm sonu çözümlü problemler}

\begin{enumerate}
  \item \textbf{Tabakalı örnekleme.} Bir fakültede sınıf düzeylerine göre öğrenci sayıları 600, 500, 400 ve 300'dür. Orantılı tabakalı yöntemle $n=180$ öğrenci seçilecektir. Her sınıftan kaç öğrenci alınmalıdır?

  \textbf{Çözüm:} Toplam $N=1800$'dür. $n_h=n(N_h/N)$ kullanılır. Sırasıyla $180(600/1800)=60$, 50, 40 ve 30 öğrenci seçilir. Sayılar 180'e toplamlanır.

  \item \textbf{Ağırlıklı oran.} Örneklemde birinci sınıftan 80 öğrencinin 32'si, dördüncü sınıftan 20 öğrencinin 12'si danışma hizmetini kullanmıştır. Evrende iki sınıfın büyüklükleri eşitse ağırlıksız örneklem oranı ile sınıf ağırlıklı oranı bulun.

  \textbf{Çözüm:} Ağırlıksız oran $(32+12)/100=0{,}44$'tür. Sınıf oranları 0,40 ve 0,60'tır. Evren payları eşit olduğundan ağırlıklı oran $0{,}5(0{,}40)+0{,}5(0{,}60)=0{,}50$ olur. Orantısız seçim dikkate alınmadığında oran küçümsenmiştir.

  \item \textbf{Yanıtlamama.} 500 kişilik örneklemde 300 kişi ankete yanıt vermiş ve bunların yüzde 40'ı yüksek kaygı bildirmiştir. Yalnız bu bilgiyle evren oranı neden güvenle tahmin edilemez?

  \textbf{Çözüm:} Yanıt oranı yüzde 60'tır. Yanıtlamayanların kaygı düzeyi yanıtlayanlardan sistematik olarak farklıysa yüzde 40 yanlı olabilir. Hatırlatma dalgaları, bilinen yardımcı değişkenlerle ağırlıklandırma ve yanıtlayan--yanıtlamayan karşılaştırması gerekir; büyük örneklem yanıtlamama yanlılığını tek başına gidermez.
\end{enumerate}

\bolumkaynaklari{\cite{cochran1977,lohr2021,chambers2003,murthy1967,cortez2008data,cortezsilva2008}}

\section*{Hatırla--Uygula--Yansıt}

\begin{enumerate}
\renewcommand{\labelenumi}{\textbf{\Alph{enumi}.}}
  \item \textbf{Hatırla:} Hedef evren, erişilebilir evren, örnekleme çerçevesi, örnekleme birimi ve analiz birimini tek örnekte ayırın.
  \item \textbf{Hesapla:} $N=2.400$ ve $n=240$ için dahil olma olasılığını ve sistematik seçim aralığını bulun.
  \item \textbf{Dağıt:} 500, 300 ve 200 kişilik üç tabakadan 100 kişilik orantılı örneklem oluşturun.
  \item \textbf{Seç:} Dağınık 80 okuldan öğrenci seçilecek bir araştırma için basit rastgele, tabakalı ve küme örneklemeyi maliyet ve hassasiyet bakımından karşılaştırın.
  \item \textbf{Tanıla:} Çevrim içi yardım arama anketinde kapsam, gönüllülük, yanıtlamama ve ölçüm yanlılığına birer örnek verin.
  \item \textbf{Eleştir:} ``10.000 kişi ankete katıldığı için sonuçlar bütün gençleri temsil eder'' savını değerlendirin.
  \item \textbf{Planla:} 400 tamamlanmış yanıt ve yüzde 80 beklenen yanıt oranı için kaç kişinin davet edilmesi gerektiğini hesaplayın.
  \item \textbf{Uygula:} Küçük bir öğrenci listesinde basit rastgele ve tabakalı örneklem oluşturup grup dağılımlarını karşılaştırın.
  \item \textbf{Etik karar:} Danışmanın kendi danışanlarını araştırmaya davet etmesindeki güç ilişkisini ve gönüllülük riskini tartışın.
  \item \textbf{Raporla:} Seçim, davet, yanıt ve dışlama sayılarını içeren kısa bir örnekleme yöntemi paragrafı yazın.
\end{enumerate}

\bolumbaglantisi{chapters/pdr/06-ornekleme.tex}